% Figure: Eckart-Young truncation error against rank k, for the
% spectral norm (= sigma_{k+1}) and the Frobenius norm (= sqrt of the
% tail energy), on a synthetic exponentially decaying spectrum.
\begin{figure}[htbp]
\centering
\begin{tikzpicture}
\begin{semilogyaxis}[
  width=13cm,
  height=8cm,
  xlabel={retained rank $k$},
  ylabel={truncation error},
  xmin=0, xmax=20,
  ymin=0.02, ymax=6,
  xtick={0,4,8,12,16,20},
  tick label style={font=\footnotesize},
  label style={font=\small},
  legend style={font=\footnotesize, at={(0.97,0.97)}, anchor=north east},
  title={Eckart-Young error: the two norms fall at different rates},
  title style={font=\bfseries},
  ymajorgrids=true,
  grid style={gray!25},
]
% Spectral-norm error at rank k equals sigma_{k+1}. Spectrum decays as
% sigma_i = 5 * 0.7^{i-1}, so error_spec(k) = 5 * 0.7^k.
\addplot[engblue, very thick, domain=0:20, samples=21] {5*0.7^x};
\addlegendentry{spectral norm $= \sigma_{k+1}$}
% Frobenius-norm error at rank k = sqrt(sum_{i>k} sigma_i^2). With the
% geometric spectrum, sum_{i>k} sigma_i^2 = 25 * 0.49^k / (1-0.49).
% error_frob(k) = sqrt(25/0.51) * 0.7^k = 7.0014 * 0.7^k.
\addplot[engred, very thick, dashed, domain=0:20, samples=21] {7.0014*0.7^x};
\addlegendentry{Frobenius norm $= \sqrt{\sum_{i>k}\sigma_i^2}$}
% mark a decision level.
\addplot[gray, thin, dotted, domain=0:20] {0.5};
\node[font=\footnotesize, gray, anchor=south west] at (axis cs:0.3,0.53)
  {tolerance $= 0.5$};
\end{semilogyaxis}
\end{tikzpicture}
\caption[Eckart-Young truncation error in two norms]{The exact error
of the best rank-$k$ approximation $\mat{A}_k$ against $k$, for a
matrix whose singular values decay geometrically as $\sigma_i =
5\,(0.7)^{i-1}$. The Eckart-Young theorem makes both curves exact,
not asymptotic: the spectral-norm error (blue) is precisely the next
singular value $\sigma_{k+1}$, and the Frobenius-norm error (red
dashed) is precisely the square root of the discarded tail energy
$\sum_{i>k}\sigma_i^2$. Both fall geometrically with the same slope
here because the spectrum is geometric, with the Frobenius error the
larger of the two at every $k$ because it accumulates the whole tail
rather than its leading term. The engineer reads off the smallest $k$
whose error clears the application tolerance (dotted): at tolerance
$0.5$ the spectral norm permits $k = 7$ while the stricter Frobenius
criterion demands $k = 8$, the one-component gap the two readings
routinely produce.}
\label{fig:vol02:ch10:eckart-young-error}
\end{figure}
